\documentclass[10pt,a4paper]{amsart}
\usepackage[margin=19mm]{geometry}
\usepackage{amsmath,amssymb,mathtools}
\usepackage[scr=rsfs,cal=euler]{mathalfa}
\usepackage{xfrac}
\usepackage[unicode,hidelinks,bookmarks=false]{hyperref}
\newcommand{\ie}{i.e.\ }
\newcommand{\cf}{cf.\ }
\newcommand{\resp}{respectively\ }
\newcommand{\Subsection}{\subsection}
\providecommand{\nobreakdash}{\nobreak\hbox{-}}
\setlength{\emergencystretch}{2em}
\multlinegap0pt
\usepackage{bm}
\usepackage{bbm}
\usepackage{dsfont}
\usepackage{xspace}
\usepackage{cancel}
\usepackage{mathtools}
\usepackage{amsthm}
\makeatletter
\@ifundefined{thm}{\newtheorem{thm}{Thm}}{}
\@ifundefined{lem}{\newtheorem{lem}[thm]{Lemma}}{}
\makeatother
\title[On spectral simplicity of the Hodge Laplacian and Curl Opera]{On spectral simplicity of the Hodge Laplacian and Curl Operator along paths of metrics}
\author{Willi Kepplinger}
\date{Source: arXiv:2211.06916v3 (CC BY 4.0)}
\begin{document}
\maketitle

\noindent\emph{Source and licence.} On spectral simplicity of the Hodge Laplacian and Curl Operator along paths of metrics. Version arXiv:2211.06916v3 (CC BY 4.0), distributed under the Creative Commons Attribution 4.0 International licence, as recorded in the arXiv licence metadata for this exact version and shown on the abstract page https://arxiv.org/abs/2211.06916v3. Full rights notice: licenses/arxiv-2211.06916-CC-BY-4.0.txt.\par\smallskip

\noindent\textit{Editorial scope.} Counted unit: the Lem whose source label is \texttt{lemma: transversality of curl operator} in the article (source0.tex lines 382--385, source label \texttt{lemma: transversality of curl operator}), delivered together with the article's own complete original proof of it (source0.tex lines 386--412, one \texttt{proof} environment of 27 lines). No mathematics is written, repaired or re-derived: statement and proof are the article's own text line for line. Required context retained uncounted: none. Every definition, display and label of those lines is retained unchanged. Omitted from this package and not redistributed: 1--381, 413--549. Nothing inside a retained excerpt is omitted or altered: every retained body is delivered exactly as the article's lines are written, so no omission receipt of any kind accompanies this package. Citations: the retained excerpts carry no citation command at all, so no bibliography entry of the article is transcribed and no reference is redistributed. Reference adaptation: none was needed and none was made; every reference inside the delivered unit resolves inside it, so no unresolved reference or citation is produced. Printed slips kept literal: no printed word, symbol, hypothesis or number of the counted statement or of its proof was altered. Layout and class: the article's own class is \texttt{amsart} with packages including amsmath, amsthm, amssymb, graphicx, biblatex, hyperref, lineno; the wrapper is the lane's pinned \texttt{amsart} editorial wrapper at 10pt on A4, which restates the theorem-like environments the retained text needs and adds the article's own macro definitions that the retained text needs (none), with extra wrapper packages (bm, bbm, dsfont, xspace, cancel, mathtools). The delivered numbering is the unit's own. Assets: no retained span contains a figure environment, a graphics-inclusion command, a PDF-page inclusion, a table or a TikZ picture; no image file is copied into this package, the package declares no asset dependency and no image file is redistributed with it. Licence: version arXiv:2211.06916v3 of this article is distributed under the Creative Commons Attribution 4.0 International licence, as recorded in the arXiv licence metadata for this exact version and shown on the abstract page https://arxiv.org/abs/2211.06916v3. A full rights notice is delivered at licenses/arxiv-2211.06916-CC-BY-4.0.txt. Characters: every retained excerpt is pure ASCII, so the delivered engine, XeLaTeX with its default Latin Modern text font, reports no missing character.
\begin{lem}
\label{lemma: transversality of curl operator}
    $Id$, $A^{\prime}[\Tilde{h}_a]$, and $A^{\prime}[v_1\odot v_2]$ span $\mathcal{L}(E(g))$ for some choice of $a$.
\end{lem}

\begin{proof}
    Setting the determinant associated to these three vectors to zero and solving for $\int_M \lVert v_1\rVert^4\,d\mu_g$ as in the preceding Lemma we get

\begin{align*}
&\int_M\lVert v_1\rVert^4\,d\mu_g=\bigg(\frac{2}{1-a}\int_M g(v_1,v_2)^2\,d\mu_g-\frac{1+a}{1-a}\int_M \lVert v_1\rVert^2 \lVert v_2\rVert ^2\,d\mu_g\bigg)+\\
&\frac{\bigg(\int_M \lVert v_1 \rVert^2 g(v_1,v_2)\,d\mu_g\bigg)^2-\bigg(\int_M \lVert v_1 \rVert^2 g(v_1,v_2)\,d\mu_g\bigg)\,\bigg(\int_M \lVert v_2 \rVert^2 g(v_1,v_2)\,d\mu_g\bigg)}{\int_M \lVert v_1\rVert^2 \lVert v_2\rVert^2\,d\mu_g}\\
& =V_a
\end{align*}
Now suppose that in fact all of these determinants vanish, so the RHS of the above equation does not depend on the choice of $a\in \mathbb{R}\setminus \{1\}$. Then in particular, $V_a=V_0$, and so we get 
\begin{align*}
    &(1-a)\bigg(2\int_M g(v_1,v_2)^2\,d\mu_g-\int_M \lVert v_1\rVert^2\lVert v_2\rVert^2\,d\mu_g\bigg)\\
    =&\bigg(2\int_M g(v_1,v_2)^2\,d\mu_g-(1+a)\int_M \lVert v_1\rVert^2\lVert v_2\rVert^2\,d\mu_g\bigg)
\end{align*}
which is equivalent to 
\begin{align*}
    \int_M \big(g(v_1,v_2)^2-\lVert v_1\rVert^2\lVert v_2\rVert^2\big)\,d\mu_g=0
\end{align*}
meaning that $v_1\parallel v_2$ for every point in $M$. We will now show that this is impossible.
For this note that the complements of the zero set of $v_1$ and $v_2$ are open and dense by the unique continuation property. This means that there exists an open set $U\subset M$ in the complement of these zero sets and a smooth function $s$ such that $v_1=s\,v_2$ on $U$. For constant $s$ we immediately get a contradiction since $\langle v_1,v_2\rangle_g=0$, so $s$ is some nonconstant function. Plugging this into the eigenform-equation for the curl operator yields
\begin{align*}
    \lambda v_1=A_g v_1=A_g (s\,v_2)=s\,A_g v_2 +\ast_g (ds\wedge v_1)=\lambda\,v_1+\ast_g (ds\wedge v_1)
\end{align*}
This means that $ds\parallel v_1$, so $v_1=f\,ds$ for some smooth function $f$. Restricting to a subset $\Tilde{U}\subset U$ on which $ds$ does not vanish, we see that this leads to the following contradiction
\begin{align*}
    0 < v_1\wedge dv_1 = f\,ds\wedge d(f\,ds)=f^2\,ds\wedge ds =0
\end{align*}
\end{proof}

\end{document}
